\documentclass[a4paper]{article}

\usepackage{graphicx}
\usepackage{amsmath,amssymb}
\usepackage{minted}
\usepackage{multirow}
\usepackage{float}
\usepackage{algorithm}
\usepackage{algpseudocode}
\usepackage{todonotes}
\usepackage{forest}
\usepackage{xstring}
\usepackage{xcolor}
\usepackage[most]{tcolorbox}
\usepackage{xparse}
\usepackage{natbib}
\usepackage{adjustbox}
\usepackage[active,tightpage]{preview}

\PreviewBorder=4pt

\usetikzlibrary{graphs,graphdrawing,arrows.meta,positioning,calc,fit,shapes.geometric}
\usegdlibrary{layered}

\input{/nd_language_ai_project/src/nd_language_ai/formal/computer/programming/tex/latex/code_clips/external_graphviz.tex}
\input{/nd_language_ai_project/src/nd_language_ai/formal/computer/programming/tex/latex/part/control_sequence/kind/semtag/semtag.tex}
\input{/nd_language_ai_project/src/nd_language_ai/formal/computer/programming/tex/latex/code_clips/seehere.tex}
\input{/nd_language_ai_project/src/nd_language_ai/formal/computer/programming/tex/latex/code_clips/deep_structure.tex}

\title{Tab Overflow Popup Hover Fix}
\date{2026-09-21}
\author{Mohammad Rahmani}

\begin{document}
   \maketitle

   \begin{preview}
      \section{Bug}
      The blue rectangle shown in the Pinned/Open overflow popup was not the
      QTabBar tab-hover paint.  It was the native popup's active item, which
      Qt keeps active after the pointer leaves that item.

      \section{Fix}
      The proxy tab bar now owns a popup-hover controller for its native
      overflow arrow buttons.  The controller styles the native menu/list
      hover in light purple and explicitly clears the popup active/current
      item on Leave, HoverLeave, hide, and window deactivation.  The existing
      native overflow mechanism and arrows remain unchanged.
   \end{preview}

   \bibliographystyle{plainnat}
   \bibliography{/home/donkarlo/Dropbox/repo/nd_shared_research/bibliography/bibliography.bib}
\end{document}
